\needspace{0.25\textheight}
\section{Инварианты: Казимир, гамильтониан, затухание}\label{sec:inv}

\subsection{Точные законы эволюции}

\begin{theorem}[Казимир]\label{th:casimir}
Вдоль решений \eqref{eq:epd} при \textup{(H1)}
\begin{equation}\label{eq:casimir}
  \frac{d}{dt}\|\Om\|^2
  = 2\,\bigl\langle \Om,\; \Om\times\bbI^{-1}\Om\bigr\rangle
    - 2\nu\,\langle\Om,\bbL\Om\rangle
  = -2\nu\,\langle\Om,\bbL\Om\rangle \;\le\; 0 .
\end{equation}
При $\nu = 0$ величина $\|\Om(t)\|^2 \equiv \|\Om_0\|^2$ тождественно:
траектория заперта на сфере $S^2_{\|\Om_0\|} \subset \so^{*}$.
\end{theorem}

\begin{proof}
Первое равенство~--- дифференцирование квадрата нормы и подстановка
\eqref{eq:epd}. Обращение в нуль перекрёстного члена~--- определитель
с двумя равными столбцами:
$\langle\Om, \Om\times\om\rangle = \det[\Om,\Om,\om] = 0$
при любых $\om$. Знак~--- положительная полуопределённость $\bbL$
из \textup{(H1)}: $\langle\Om,\bbL\Om\rangle \ge 0$.
\end{proof}

\begin{theorem}[гамильтониан]\label{th:ham}
Вдоль решений \eqref{eq:epd}
\begin{equation}\label{eq:ham}
  \frac{d\mathcal H}{dt}
  = \bigl\langle \nabla\mathcal H,\dot\Om\bigr\rangle
  = \underbrace{\bigl\langle \bbI^{-1}\Om,\;\Om\times\bbI^{-1}\Om\bigr\rangle}_{=0}
  - \nu\,\bigl\langle \bbI^{-1}\Om,\;\bbL\Om\bigr\rangle
  = -\nu\,\Om^{\mathsf T}\bbI^{-1}\bbL\,\Om .
\end{equation}
Если дополнительно выполнена \textup{(H2)}, то
\begin{equation}\label{eq:ham-diag}
  \frac{d\mathcal H}{dt} = -\nu\sum_{k=1}^{3}\frac{L_{kk}}{I_k}\,\Omega_k^2 \;\le\; 0 ,
\end{equation}
и $\mathcal H$~--- функция Ляпунова строгого сертификата устойчивости.
\end{theorem}

\begin{proof}
Равенство \eqref{eq:ham}~--- правило цепочки с
$\nabla\mathcal H = \bbI^{-1}\Om$ и обнуление перекрёстного члена как в
теореме~\ref{th:casimir}. При \textup{(H2)} операторы $\bbI^{-1}$ и $\bbL$
диагональны в одном базисе, произведение $\bbI^{-1}\bbL$ симметрично и
положительно полуопределено, откуда \eqref{eq:ham-diag}.
\end{proof}

\begin{remark}[коммутационная тонкость]\label{rem:commutator}
Гипотеза \textup{(H2)} в теореме~\ref{th:ham} существенна и не сводится
к \textup{(H1)}: без неё квадратичная форма
$\Om \mapsto \Om^{\mathsf T}\bbI^{-1}\bbL\Om$ определяется
симметризацией $\tfrac12(\bbI^{-1}\bbL + \bbL\bbI^{-1})$, которая может
быть неопределённой даже при $\bbL \succ 0$. Пример:
\[
  \bbI^{-1} = \diag(10,1,1),\qquad
  \bbL = \begin{pmatrix} 1 & 1 & 0\\ 1 & 1{,}5 & 0\\ 0 & 0 & 1 \end{pmatrix},
  \qquad \spec \bbL = \{0{,}219;\;1;\;2{,}281\} \subset (0,\infty),
\]
\[
  \operatorname{sym}(\bbI^{-1}\bbL) = \begin{pmatrix} 10 & 5{,}5 & 0\\ 5{,}5 & 1{,}5 & 0\\ 0 & 0 & 1 \end{pmatrix},
  \qquad \det\begin{pmatrix} 10 & 5{,}5\\ 5{,}5 & 1{,}5 \end{pmatrix} = -15{,}25 < 0 .
\]
В точке $\Om_0 = (-0{,}5;\,1;\,0)$ имеем
$\Om_0^{\mathsf T}\bbI^{-1}\bbL\Om_0 = -1{,}5$, т.е.
$\dot{\mathcal H}(0) = +1{,}5\,\nu > 0$: \emph{гамильтониан возрастает,
тогда как Казимир по-прежнему строго убывает} (теорема~\ref{th:casimir}
не требует \textup{(H2)}). Численная иллюстрация~--- проверка E1 в
\S\ref{sec:numerics}: за $t \le 4$ рост $\mathcal H$ на $1{,}02\%$ при
строго убывающем $\|\Om\|^2$. Таким образом, «сертификат Ляпунова»
гарантирован только при \textup{(H2)}; энергетическая монотонность~---
свойство диагонализации, а не диссипации самой по себе.
\end{remark}

\subsection{Скорость затухания и предельное множество}

\begin{corollary}[экспоненциальное затухание]\label{cor:decay}
Если $\lambda_{\min} := \lambda_{\min}(\bbL) > 0$, то
\begin{equation}\label{eq:decay}
  \|\Om(t)\|^2 \le \|\Om_0\|^2\,e^{-2\nu\lambda_{\min} t},
  \qquad t \ge 0 .
\end{equation}
\end{corollary}

\begin{proof}
Из \eqref{eq:casimir} и неравенства Рэлея
$\langle\Om,\bbL\Om\rangle \ge \lambda_{\min}\|\Om\|^2$:
$\dot C \le -2\nu\lambda_{\min}C$ для $C = \|\Om\|^2$. Тогда
$\varphi(t) := C(t)e^{2\nu\lambda_{\min}t}$ удовлетворяет
$\dot\varphi = e^{2\nu\lambda_{\min}t}\bigl(\dot C + 2\nu\lambda_{\min}C\bigr)
\le 0$, т.е. $\varphi$ невозрастает и $\varphi(t) \le \varphi(0)$.
\end{proof}

\begin{theorem}[Ла--Салль]\label{th:lasalle}
Пусть $\omega$-предельное множество траектории \eqref{eq:epd}
непусто и компактно (это так при $\nu > 0$). Тогда оно содержится в
наибольшем инвариантном подмножестве множества
$M_0 = \{\Om : \langle\Om,\bbL\Om\rangle = 0\} = \ker\bbL$. В частности:
\begin{enumerate}[label=\textup{(\alph*)}, leftmargin=2.2em, itemsep=1pt]
  \item при $\bbL \succ 0$: $\Om(t) \to 0$ при $t\to+\infty$ для любого
        начального условия;
  \item при $\nu > 0$, диагональных $\bbL$, $\bbI$ и одномерном
        $\ker\bbL = \operatorname{span}\{e_k\}$: $\Om(t)$ сходится к
        точке равновесия $a\,e_k$, $a \in \R$.
\end{enumerate}
\end{theorem}

\begin{proof}
По теореме~\ref{th:casimir} $C = \|\Om\|^2$ невозрастает и ограничена
снизу, значит существует $C_\infty = \lim C(t)$. На $\omega$-предельном
множестве $C \equiv C_\infty$ и $\dot C = 0$; по принципу
Ла--Салля \cite{lasalle} $\omega$-предельное множество инвариантно и
лежит в $M_0 = \ker\bbL$. (a): $\ker\bbL = \{0\}$, и единственная
инвариантная точка в $\{0\}$~--- само $0$ (равновесие). (b): орбита,
целиком лежащая в $\operatorname{span}\{e_k\}$, имеет вид
$\Om(t) = a(t)e_k$; подстановка в \eqref{eq:epd} даёт
$\dot a\,e_k = a^2\,e_k\times(e_k/I_k) = 0$, откуда $a$ постоянна:
инвариантные точки в $\ker\bbL$~--- в точности равновесия $ae_k$.
\end{proof}

\subsection{Топология орбит при \texorpdfstring{$\nu=0$}{nu=0}}

\begin{proposition}[периодичность]\label{prop:periodic}
Пусть $\nu = 0$, моменты $I_1, I_2, I_3$ попарно различны и
$\Om_0$ не лежит на осях $\{ \pm a e_k \}$. Тогда орбита
$\{\Om(t)\}$~--- замкнутая гладкая кривая (окружность Лиувилля),
и движение по ней периодично.
\end{proposition}

\begin{proof}
\emph{Шаг 1: равновесия.} При $\nu = 0$ равновесия удовлетворяют
$\Om\times\bbI^{-1}\Om = 0$, т.е. $\Om \parallel \bbI^{-1}\Om$~---
$\Om$ собственный вектор $\bbI^{-1}$, значит $\Om \in \bigcup_k
\operatorname{span}\{e_k\}$; на сфере $S^2_R$ это шесть точек
$\pm Re_k$, $R = \|\Om_0\|$.

\emph{Шаг 2: критические точки $\mathcal H|_{S^2_R}$.} Множитель
Лагранжа для условного экстремума $\mathcal H$ на сфере даёт
$\Omega_k(1/I_k - \lambda) = 0$ для всех $k$; нетривиальные решения~---
снова $\pm Re_k$. В этих шести точках $\mathcal H|_{S^2_R}$ принимает
три попарно различных значения $\tfrac{R^2}{2I_k}$, и все они
невырождены (попарное различие $I_k$).

\emph{Шаг 3: структура регулярных уровней.} Пусть
$h \notin \{\tfrac{R^2}{2I_k}\}$ и
$L_h = \{\Om \in S^2_R : \mathcal H(\Om) = h\}$. На $L_h$ градиент
$\mathcal H|_{S^2_R}$ не обращается в нуль, поэтому по теореме о
неявной функции $L_h$~--- гладкое компактное одномерное многообразие
без края; по классификации компактных одномерных многообразий каждая
его компонента диффеоморфна окружности. По теореме~\ref{th:casimir} и
предложению~\ref{prop:poisson} ($\nu = 0$: гамильтонов поток) обе
функции $C$ и $\mathcal H$ сохраняются, значит каждая компонента
$L_h$ инвариантна относительно потока.

\emph{Шаг 4: периодичность.} На $L_h$ векторное поле \eqref{eq:ep} не
обращается в нуль (нули~--- в точности равновесия шага~1, лежащие на
критических уровнях). Параметризуем компоненту углом $\theta$; тогда
$\dot\theta = a(\theta)$, где $a$ непрерывна и без нулей на
компактной окружности, значит $|a| \ge a_0 > 0$ и время полного
обхода
\[
  T = \oint \frac{d\theta}{|a(\theta)|} \le \frac{2\pi}{a_0} < \infty .
\]
Следовательно каждая орбита~--- вся компонента $L_h$, проходимая за
конечное время $T$: движение периодично.
\end{proof}

\begin{proposition}[устойчивость равновесий]\label{prop:stability}
Пусть $\nu = 0$, $I_1 > I_2 > I_3$. На сфере $S^2_R$ точки $\pm Re_1$
и $\pm Re_3$ устойчивы по Ляпунову, точки $\pm Re_2$ неустойчивы:
линеаризация в $Re_2$ имеет действительные собственные значения
\[
  \lambda_{\pm} = \pm R\sqrt{\Bigl(\tfrac1{I_1}-\tfrac1{I_2}\Bigr)
                           \Bigl(\tfrac1{I_2}-\tfrac1{I_3}\Bigr)} \ne 0 .
\]
\end{proposition}

\begin{proof}
\emph{Устойчивость $\pm Re_1$, $\pm Re_3$.} На $S^2_R$ вблизи $Re_1$
($I_1$~--- максимум)
\[
  \mathcal H - \mathcal H(Re_1)
  = \tfrac12\sum_{k=2}^{3}\Bigl(\tfrac1{I_k}-\tfrac1{I_1}\Bigr)\Omega_k^2
  \ge c\,\bigl(\Omega_2^2+\Omega_3^2\bigr), \qquad c > 0,
\]
т.е. $\mathcal H|_{S^2_R}$ имеет строгий локальный минимум в $Re_1$;
сохранение $\mathcal H$ и тот факт, что малые уровни $\mathcal H|_{S^2_R}$
суть малые замкнутые кривые вокруг $Re_1$ (следствие
предложения~\ref{prop:periodic}), дают устойчивость по Ляпунову;
для $Re_3$ ($I_3$~--- минимум) аналогично с максимумом $-\mathcal H$.

\emph{Неустойчивость $Re_2$.} Положим $\Om = Re_2 + \xi$,
$\|\xi\|$ мало. Линеаризация \eqref{eq:ep}:
\[
  \dot\xi = Re_2\times\bbI^{-1}\xi + \xi\times\bbI^{-1}(Re_2)
  = R\Bigl(\tfrac{\xi_3}{I_3},\,0,\,-\tfrac{\xi_1}{I_1}\Bigr)
    + \Bigl(-\tfrac{R}{I_2}\xi_3,\,0,\,\tfrac{R}{I_2}\xi_1\Bigr),
\]
и после приведения
\[
  \dot\xi_1 = -R\Bigl(\tfrac1{I_2}-\tfrac1{I_3}\Bigr)\xi_3,
  \qquad
  \dot\xi_3 = R\Bigl(\tfrac1{I_2}-\tfrac1{I_1}\Bigr)\xi_1,
  \qquad \dot\xi_2 = 0 .
\]
Матрица $\begin{psmallmatrix}0 & -a\\ b & 0\end{psmallmatrix}$,
$a = R(\tfrac1{I_2}-\tfrac1{I_3}) < 0$,
$b = R(\tfrac1{I_2}-\tfrac1{I_1}) > 0$, имеет
$\lambda_{\pm}^2 = -ab > 0$, т.е. $\lambda_+ > 0$~--- по теореме об
неустойчивости по первому приближению \cite[гл.~VII]{hartman}
равновесие неустойчиво (эффект «теннисной ракетки»/Джанибекова).
\end{proof}
